%!TEX root = ../Thesis.tex
\chapter{Introduction}

\section{Background}

When the British archaeologist Arthur Evans in the beginning of the 20th century went to excavate Knossos on Crete, he found a huge amount of clay tablets written in a previously unknown linear script.
Some documents were clearly of a newer script, which he called \textit{Linear B}, while others were of an older one, which he named \textit{Linear A}.
Some fifty years later, the British amateur archaeologist Michael Ventris finally succeeded in deciphering the newer script Linear B, which he proved to be representing Mycenaean Greek, the earliest attested form of Greek.
But the same hypothesis did not hold for the older script Linear A, whose encoded language remained unknown, still yet to be deciphered.

A key obstacle to its decipherment has been the small size of its corpus, lately consisting of only around 2,500 inscriptions (about 1,400 if only including those well-preserved enough to be readable) - significantly smaller than that of its younger counterpart that now spans more than 4,500 inscriptions (which are comparatively longer and better preserved).\footcite{Salgarella2020}\footcite[p.~13,46]{Salgarella2025} 
But new inscriptions are still being discovered, with the corpus most recently extended in 2025 by over 100 new documents.\footcite{RILA-S1}\footcite{CNRLinearACorpus2025}

Another challenge has been the lack of digital representations of these inscriptions, in a form suitable for statistical analysis.
While photographs with accompanied by expert drawings of nearly the entire corpus - \acf{GORILA} - are available online through \acf{Cefael} (see \textcite{EtudesCretoises21}), these are only available in print form.
This motivated the development of the palæographical\footnote{The term "palæographical" refers to the study of ancient writing systems.} database \textit{The Signs of Linear A} (SigLA), which seeks to "[develop] a systematic, exhaustive and user-friendly open access database of all Linear A inscriptions [...] in order to carry out statistical and palæographic analyses within the epigraphic corpus" \parencite{SalgarellaCastellanSigLA2020}.
However, the last addition to SigLA was in 2021, and the project has so far managed to document just 772 of all currently known inscriptions.\footcite{SigLAAboutPage}\footcite{SigLASearchDocuments}

\subsection{A problem of manual digitization}

The drawings currently stored on SigLA were made by Ester Salgarella using the painting tool Krita.
For each document, she manually created her own annotated drawings based on the \acs{GORILA} originals stored by \acs{Cefael} as raw tablet images accompanied by expert drawings.
From these, she proceeded to create layered Krita files with metadata for each sign, which could then at last be exported as JSON files to SigLA.\footcite{SalgarellaCastellanSigLA2020}
According to her (personal communication), this has been a highly time-consuming undertaking. 
And the necessity of this kind of labor-intensive work has arguably been the biggest reason to why SigLA still does not contain the entire available corpus.

Due to the potential violation of \acf{IPR}, it is also of particular interest, whether this process can begin from the digital source images themselves, before relying on already completed but undigitized expert drawings as input.
The aim of this project was therefore to investigate whether parts of the process can be automated by working on the raw tablet images from \acs{GORILA}.
More specifically, whether some kind of semi-automatic segmentation tool (given just the raw tablet images as input) can accelerate the workflow, while maintaining high-quality outputs that are still compatible with SigLA's database.

\subsection{A problem of small data}
The small size of the Linear A corpus and the irregularities in quality of its documents is a huge problem for the application of deep learning approaches, which typically require large amounts of data to perform well.
Consequently, this project instead sought to make use of generalized \acf{ML} and \acf{CV}.
In particular Meta's \acf{SOTA} interactive segmentation model, the \acf{SAM} \parencite{sam}, which is pre-trained on a massive and diverse dataset - \acf{SA-1B}, of over 1 billion masks.
\acs{SAM}'s generalization and interactivity (taking user-placed prompts that guide the segmentation) potentially allows for semi-automatic annotation of Linear A inscriptions without the need for large amounts of training data, which was unfortunately not available in this context.
It was therefore of particular interest to map the capabilities and limitations of \acs{SAM} within this domain.


\section{Related work}
Within the last decade, the field of document annotation has seen significant advancements through the use of machine learning.
Take for example the EU-funded generic \acs{CNN}-based framework for historical document processing \textit{dhSegment} \parencite{dhsegment}.
Or the more recent work on segmentation of stone inscriptions \parencite{DeepLearningSegmentation}, where \acs{SAM} and other state-of-the-art methods are beaten as benchmarks by their stacked U-Net and \acs{GAN} architecture.
Or even the \acs{YOLO} and Roboflow based approach to segmentation of Palmyrene Aramaic inscriptions \parencite{YOLOSegmentation} that augment\footnote{The term \textit{augmenting} refers to the process of increasing the size of a dataset.} their dataset with a \acf{GAN} to increase performance.

Nevertheless, these methods all depend on deep learning, making them inconvenient for domains of small data, which instead seem to require a generalized and interactive approach such as \acs{SAM}.

\subsection{SAM and its applications}
Ever since its release in 2023, \acs{SAM} has been extensively evaluated across a wide range of \acs{CV} tasks to investigate whether it can really "segment anything".
Papers have evaluated its performance within niche domains, such as medical images \parencite{MedSAM}, camouflaged objects \parencite{CamSAM}, transparent objects \parencite{TransparentSAM}, hierarchical texts \parencite{HiSAM}, and even oracle bone inscriptions \parencite{OBIFloSAM}. 
Most of these papers found that \acs{SAM} still tends to struggle with such domains, despite its impressive zero-shot\footnote{The term \textit{zeroshot} refers to model performance on hitherto unseen data.} generalization.
Meanwhile, some of them also found that this performance can be significantly improved through architectural extensions and domain-specific fine-tuning, which unfortunately requires deep learning and a lot of data.
However, equally, it has been shown that \acs{SAM}'s performance can significantly improve through effective prompt engineering alone \parencite{promptSAM}, requiring no deep learning at all.
And within the domain of 3D medical images, a training-free few-shot adaptation of \acs{SAM}2 \parencite[see][]{sam2} has been recently demonstrated to surpass supervised methods like U-Net, and even the best fined-tuned versions of \acs{SAM} \parencite{fatesam}.
Thus, it seemed reasonable to test out the limits of \acs{SAM}'s generalization and interactivity within the domain of Linear A document images, to investigate whether it could be used to accelerate the annotation workflow of Ester Salgarella.

\section{Problem statement}
As previously stated, the aim of this project was to investigate whether parts of the process of creating annotated drawings like the ones in SigLA can be automated by working on the raw tablet images from \acs{GORILA}.
More specifically, whether some kind of semi-automatic segmentation tool (given just the raw tablet images as input) using interactive and generalized \acs{CV} - specifically the \acf{SAM} - can accelerate the workflow, while maintaining high-quality outputs that are still compatible with a database like SigLA.
But what would be the relative segmentation quality of such a tool? How interaction efficient would it be exactly? And by how much could it improve the efficiency of the workflow?
To concretize the research, the problem statement was decomposed into three \acfp{RQ} - each accompanied by their respective operationalization (see the following subsection for details).

\subsection{Research questions}
\label{research_questions}
\begin{enumerate}
\item \textbf{\acs{RQ}1 (Segmentation quality)}: 
\textit{How well does a \acs{SAM}-based segmentation approach perform quantitatively in terms of segmentation accuracy, compared to simple classical image processing baselines when applied to Linear A document images?}

\textbf{Operationalization}: \acs{SAM}-based masks were compared to classical image-processing baselines using overlap metrics (i.e. \acs{Dice}) against the ground truth segmentations derived from SigLA's Krita files using a self-made dataset of 60 documents.
\item \textbf{\acs{RQ}2 (Interaction efficiency)}: 
\textit{To which extent can user interaction be reduced by a \acs{SAM}-based semi-automatic segmentation workflow, while still maintaining high-quality outputs?}

\textbf{Operationalization}: Two self-developed segmentation tools - one SAM-based, one classical - were evaluated in a case study of 3 documents. The time spent in each tool by Ester Salgarella from image import to mask export was compared against the segmentation quality (i.e. \acs{Dice}) relative to ground truth (SigLA's Krita files).
\item \textbf{\acs{RQ}3 (Workflow efficiency)}: 
\textit{How much reduction in annotation time can be achieved using the proposed tool compared to fully manual digitization, under controlled experimental conditions?}

\textbf{Operationalization}: The time it took Ester Salgarella to annotate manually compared to being assisted by two self-developed segmentation tools - one based on SAM, and the other based on a classical baseline - was measured in a case study of 3 documents.
\end{enumerate}

\section{Impact -- innovation and application}

The work of this thesis presents itself as a first proof-of-concept (and a mapping of the difficulties involved) in using \acs{SAM} and modern \acs{CV} to accelerate the annotation of Linear A documents.
Its main contribution is a prototype of a semi-automatic annotation tool that extracts engravings from raw tablet images, along with an evaluation framework that can be used to assess the performance of that tool and any potential future derivatives.
Such a tool could potentially also be applied to Linear B, as it is the whole field of Aegean studies that seems to suffer from severe under-digitization of the primary evidence (which is a huge problem for the scientific investigations of scholars).
Hopefully, the work will inspire further research and development within the field of digital epigraphy.
Particularly in the application of generalized and interactive \acs{CV} approaches such as \acs{SAM} on Linear A, and perhaps even in the broader context of use on small corpus ancient writing systems, where the lack of big data is ill fit for approaches using deep learning.


\section{Overview of the thesis}
The rest of this thesis is structured into four main chapters (Data, Methodology, Results and Discussion, Conclusion), all split up into sections concerning each of the three research questions.
